\section{Entrega del Producto}

El Kit Pagos Colombia no tiene un cliente externo en el sentido comercial del término. Los destinatarios de la entrega del producto son dos: el Ph.D. Jaime Andrés Pavlich Mariscal en su rol de director del trabajo de grado, quien recibe los artefactos intermedios y aprueba formalmente los entregables clave a lo largo del semestre; y la Facultad de Ingeniería de la Pontificia Universidad Javeriana, que recibe el producto completo al cierre del semestre como parte de los procedimientos académicos del trabajo de grado. Esta distinción implica que el proceso de entrega tiene dos momentos: entregas parciales al director asociadas a los hitos del ciclo de vida DSR, y una entrega final consolidada que cierra el proyecto.

Las entregas parciales ocurren al alcanzar los hitos H3, H4 y H5 definidos en la sección 3.1.3. En cada uno de estos puntos, el Coordinador del equipo elabora y remite al director un reporte de progreso que incluye el estado de los artefactos completados hasta ese momento, junto con evidencias de que cumplen los criterios del Plan de Aceptación correspondiente (sección 3.3). El director valida el avance y emite retroalimentación antes de que el equipo proceda a la siguiente fase.

La entrega final del producto se concentra en el hito H6 y comprende cinco artefactos, conforme a lo establecido en la sección 1.4. La tabla a continuación detalla cada artefacto, su canal de entrega, el responsable y el momento en que se ejecuta la acción de entrega.

\begin{center}
\renewcommand{\arraystretch}{1.2}
\begin{longtable}{|P{0.194}|P{0.192}|P{0.192}|P{0.192}|P{0.194}|}
  \hline
  \textbf{Artefacto} & \textbf{Canal de entrega} & \textbf{Responsable de la entrega} & \textbf{Momento} & \textbf{Criterio de aceptación} \\
  \hline
  \endfirsthead
  \hline
  \textbf{Artefacto} & \textbf{Canal de entrega} & \textbf{Responsable de la entrega} & \textbf{Momento} & \textbf{Criterio de aceptación} \\
  \hline
  \endhead
  \caption{Entregables.}
  \label{tab:entregables-producto}\\
  \endlastfoot
  Prototipo funcional del SDK con interfaz unificada & Publicación del paquete en npm (versión estable 1.0.0) & Líder de arquitectura y núcleo SDK & H6 & Sección 3.3.1 \\
  \hline
  API de simulación desplegada & URL pública en Railway o Render, accesible desde internet & Líder de integración & H6 & Sección 3.3.2 \\
  \hline
  Documentación centralizada de datos de prueba & Repositorio público en GitHub (directorio /docs/sandbox-data) & Líder de integración & H6 & Sección 3.3.3 \\
  \hline
  Informe comparativo de métricas de evaluación (WMC, CBO, RFC) & Capítulo de evaluación dentro del documento final de grado & Líder de evaluación y documentación & H6 & Sección 3.3.4 \\
  \hline
  Documento final del trabajo de grado y repositorio público & Entrega al director y a la Facultad por los canales institucionales; repositorio público en GitHub con documentación TypeDoc & Líder de evaluación y documentación & H6 & Sección 3.3.5 \\
  \hline
\end{longtable}
\end{center}

Ningún artefacto se considera formalmente entregado si no cumple previamente los criterios de aceptación descritos en la sección 3.3. Es responsabilidad del integrante que lidera cada frente verificar el cumplimiento de dichos criterios antes de ejecutar la acción de entrega. El director aprueba la entrega del documento final y el Coordinador del equipo verifica que el repositorio público esté accesible y documentado con suficiente detalle para que un desarrollador externo pueda integrar el kit sin contactar a los autores, conforme al criterio de aceptación de la sección 3.3.5. Las acciones de archivo y cierre formal del proyecto están descritas en la sección 5.3.2.
